\subsection*{Idempotence}
\label{subsec:idempotence}

\begin{definition}[Idempotent Operation]
\label{def:idempotent-operation}
Let $\star$ be a binary operation on $A$.
The operation $\star$ is \textbf{idempotent} if
\[
\forall a \in A,\quad a \star a = a.
\]
\end{definition}

\begin{remark*}[Standard quantified statement]
\[
\operatorname{IdempotentOperation}(\star,A)
\Longleftrightarrow
\operatorname{BinaryOperation}(\star,A)
\land
\forall a\in A,\; a\star a=a.
\]
\end{remark*}

\begin{remark*}[Definition predicate reading]
$\star$ is idempotent exactly when repeating the same input does not change it.
\end{remark*}

\begin{remark*}[Negated quantified statement]
\[
\neg\operatorname{IdempotentOperation}(\star,A)
\Longleftrightarrow
\exists a\in A,\; a\star a\neq a.
\]
\end{remark*}

\begin{remark*}[Interpretation]
Idempotence says that self-combination stabilizes immediately.
\end{remark*}

\begin{dependencies}
\begin{itemize}
  \item \hyperref[def:binary-operation]{Binary Operation}
\end{itemize}
\end{dependencies}

\begin{remark*}[Regeneration Note]
The prior working notes did not contain a clearly isolated idempotence
discussion in the legacy Algebraic Structures chapter.
\end{remark*}
